\documentclass[11pt,a4paper]{article}
\usepackage[utf8]{inputenc}
\usepackage{amsmath,amssymb,amsthm}
\usepackage{geometry}
\geometry{margin=1.15cm}
\usepackage{setspace}
\setstretch{1.25}
\usepackage{booktabs}
\usepackage{enumitem}
\usepackage{hyperref}
\hypersetup{colorlinks=true,linkcolor=blue,citecolor=blue,urlcolor=blue}

\newtheorem{theorem}{Theorem}
\newtheorem{proposition}{Proposition}
\newtheorem{lemma}{Lemma}
\newtheorem{corollary}{Corollary}
\theoremstyle{definition}
\newtheorem{postulate}{Postulate}
\newtheorem{definition}{Definition}
\theoremstyle{remark}
\newtheorem{remark}{Remark}

\newcommand{\ag}{a_\gamma}
\newcommand{\CQ}{C_Q}
\newcommand{\Ok}{\Omega_K}
\newcommand{\OL}{\Omega_\Lambda}
\newcommand{\Om}{\Omega_m}
\newcommand{\Mpl}{M_{\mathrm{Pl}}}
\newcommand{\avg}[1]{\langle #1\rangle}

\title{A Parameter-Free Relation Between Spatial Curvature\\
and the Cosmological Constant,\\
Derived From the Conformal Anomaly of the Photon}
\author{Lukas R\"ohl}
\date{\today}

\begin{document}
\maketitle

\begin{abstract}
This is a self-contained account written for a mathematical reader who is \emph{not} assumed to
know quantum field theory in curved space. Every object is defined and every step is derived from
first principles as far as is possible; the one place where a derivation reaches a genuine
physical hypothesis is isolated and labelled as such. The central result is a single dimensionless
number relating the spatial-curvature density $\Ok$ and the dark-energy density $\OL$ of the
universe,
\[
\frac{\lvert\Ok\rvert}{\OL}=\frac{\ag}{8\pi^2}=\frac{31}{1440\pi^2}\approx 2.181\times10^{-3},
\]
where $\ag=31/180$ is the type-A conformal-anomaly coefficient of the free electromagnetic field.
We derive $\ag$ from the Gilkey heat-kernel computation (\S\ref{sec:heat}); the topological
prefactor from Gauss--Bonnet--Chern (\S\ref{sec:gbc}); the curvature charge $\CQ=8\pi^2$ from an
elementary integral (\S\ref{sec:Q}); the spectral backbone $R1$--$R12$, including two independent
evaluations of the closed-manifold $\zeta$-value (\S\ref{sec:spectral}); the suppression of the
naive vacuum energy by sequestering (\S\ref{sec:seq}); the lift to the cosmological scale by
coherent spectral accumulation (\S\ref{sec:lift}); the integrated invariant $4\ag$ by two
independent routes (\S\ref{sec:invariant}); the sign of the curvature from a variational principle
on the hemisphere (\S\ref{sec:cap}); and the uniqueness of the photon among candidate fields
(\S\ref{sec:uniq}). The Friedmann translation (\S\ref{sec:friedmann}) yields $\OL=4\ag=0.6889$,
$\Ok\simeq+1.5\times10^{-3}$, $w=-1$. The single physical hypothesis is stated in
\S\ref{sec:postulate}; \S\ref{sec:ledger} tabulates the logical status of every claim. An
accompanying script (\texttt{csg\_kp\_falsification.py}) recomputes every number from its inputs.
\end{abstract}

\section*{Prologue --- The Curvature that Light Betrays}
\textit{A short, non-technical preface for the general reader.}

\medskip
No one has ever seen space itself. We see the things \emph{in} space --- stars, galaxies, the
light that reaches us from them. But the shape of space we can only infer. Is it flat like a
table? Or does it curve, the way the surface of the Earth curves, except that we sit in the middle
of it and notice nothing?

This harmless question leads quickly into an embarrassment. Our best theory says the empty vacuum
ought to carry an energy --- one too large by a factor with a hundred and twenty-two zeros. It is
the worst prediction physics has ever made. For over a hundred years no one has been able to
repair it.

This work does not repair it. It follows instead a faint trail, and the trail lies in the light.

In flat, empty space, light is indifferent to every scale; it knows neither large nor small. But
curve the space and that indifference is lost --- quantum theory leaves behind a tiny residue. The
beautiful part: this residue is not a matter of choice. It is a fixed number, as inescapable as the
$\pi$ in the circumference of a circle. For light it is thirty-one one-hundred-eightieths.

And now the wager of this work: the same small number that light betrays in curved space also sets
how strongly the whole universe curves. Work it out, and a prediction emerges --- about two
thousandths, and space curves not inward but a tiny bit outward. No dial to turn. One number,
finished.

There is a catch, and the work does not conceal it. The number itself --- the bare ratio between
how much the universe curves and how much dark energy fills it --- follows as a theorem; it owes
nothing to choice. What rests on a single assumption is the bridge: that this timeless number is
carried faithfully into the universe we actually inhabit, and fixes there the absolute amount of
dark energy. That assumption is not proven; it is contested; it is placed in the open, and the
empirical claim rests visibly upon it. That is the whole shape of the work: one number (a theorem),
one assumption (the bridge), one test.

For the best thing about a prediction is that one can fail at it. In the coming years two great sky
surveys, DESI and Euclid, will measure the curvature of space more precisely than ever before ---
exactly that number. If it does not hold, the idea is dead. Not patched, not reinterpreted: dead.
This readiness to die is no flaw. It is the only thing that makes an idea honest.

What begins here, then, is not a finished answer but a cleanly placed wager. If the sky carries the
number, an astonishingly simple connection will have proven true --- between the smallest thing, a
flaw in light, and the largest, the shape of the cosmos. And if not, we will know soon. That is
almost as much worth having.


\tableofcontents

\section{What a resolution must deliver}\label{sec:problem}
The cosmological constant problem has two faces. The \emph{hierarchy} face asks why the observed
vacuum energy is not of order $\Mpl^4$, a naive expectation that overshoots by a factor
$\sim10^{122}$. The \emph{value} face asks why the dark-energy fraction is the particular number
$\OL\approx0.69$, neither $0$ nor $1$. A complete resolution should (i) explain the suppression
without fine-tuning, (ii) produce the value parameter-free, and (iii) make a falsifiable
prediction distinguishing it from a bare constant.

We do \emph{not} attempt a quantum theory of gravity. We pose a narrower and sharper question:
\emph{is there a pure number, computable from established field content alone, that the
curvature-to-dark-energy ratio of the universe is forced to equal?} The body of this paper answers
yes and derives that number. One physical hypothesis is required to convert the proven Euclidean
ratio into the present-epoch absolute value; it is flagged explicitly in \S\ref{sec:postulate} and
is the proper target of any attempt at falsification.

\paragraph{Roadmap in plain language.} A field that is classically blind to the overall scale of
distances (``conformally invariant'') generically \emph{loses} that blindness on quantization. The
residue is a geometric, number-valued quantity, the \emph{conformal (trace) anomaly}. For the
photon it is controlled by one rational coefficient, $\ag$. On the round four-sphere $S^4$ --- the
Euclidean form of de~Sitter space --- this anomaly integrates to a topological invariant. We
compute that invariant, divide by a second, elementary geometric invariant of the hemisphere, and
obtain a pure number. That number is the claimed curvature-to-dark-energy ratio.

\section{The trace anomaly from first principles}\label{sec:anomaly}
\begin{definition}[Weyl rescaling, conformal invariance]
A \emph{Weyl rescaling} of a metric is $g_{\mu\nu}(x)\mapsto e^{2\sigma(x)}g_{\mu\nu}(x)$ for a
smooth function $\sigma$. An action $S[g,\phi]$ is \emph{conformally invariant} if it is unchanged
under a Weyl rescaling accompanied by the appropriate rescaling of the fields $\phi$. The free
Maxwell action $S=-\tfrac14\int\sqrt g\,F_{\mu\nu}F^{\mu\nu}$ in four dimensions is conformally
invariant, and consequently its classical stress tensor is traceless: $T^\mu{}_\mu=0$.
\end{definition}

\begin{definition}[Trace anomaly and the type-A coefficient]
The path-integral measure of the quantized field is not Weyl-invariant; the regularized one-loop
effective action $\Gamma[g]$ acquires a dependence on the scale $\sigma$. The trace of the
quantum stress tensor is then nonzero,
\begin{equation}\label{eq:anomaly}
\avg{T^\mu{}_\mu}=\frac{1}{16\pi^2}\bigl(c\,W^2-\ag\,E_4\bigr)+\nabla_\mu J^\mu,
\end{equation}
where $W_{\mu\nu\rho\sigma}$ is the Weyl tensor,
$E_4=R_{\mu\nu\rho\sigma}R^{\mu\nu\rho\sigma}-4R_{\mu\nu}R^{\mu\nu}+R^2$ is the Euler (Pfaffian)
density, and $J^\mu$ a scheme-dependent current. The coefficient $c$ multiplies the conformally
invariant $W^2$ (``type-B'') and $\ag$ multiplies the topological $E_4$ (``type-A''). Both are
pure numbers determined by the field content \cite{DeserSchwimmer,Duff1994}.
\end{definition}

\begin{remark}\label{rem:why}
Two structural facts drive everything below. \emph{(a)} On a conformally flat manifold (such as
$S^4$) the Weyl tensor vanishes, $W=0$, so only the type-A term contributes. \emph{(b)} The
integral $\int_M E_4$ over a closed manifold is a topological invariant (Theorem~\ref{thm:gbc}).
Hence on $S^4$ the integrated anomaly is a pure number times $\ag$, free of any local scheme
ambiguity.
\end{remark}

\paragraph{Why only the photon.} A field of mass $m$ contributes to the low-energy vacuum response
through terms suppressed by powers of $m^2/\Mpl^2$ and decouples in the infrared. The photon is
exactly massless --- it is the unbroken $U(1)$ gauge boson of electromagnetism --- hence exactly
conformally coupled, and its type-A anomaly is the clean, unsuppressed contribution at the order we
need. It is the sole field entering the construction; the uniqueness of this choice is established
in \S\ref{sec:uniq}.

\section{The coefficient $\ag=31/180$ from the heat kernel}\label{sec:heat}
\begin{theorem}[Photon type-A coefficient]\label{thm:agamma}
The type-A conformal-anomaly coefficient of the free Maxwell field is $\ag=\dfrac{31}{180}$.
\end{theorem}
\begin{proof}
For a second-order elliptic operator $D=-\Box+E$ acting on sections of a vector bundle over a
four-manifold, the heat kernel has the DeWitt--Schwinger short-time expansion
$K(x,x;s)=(4\pi s)^{-2}\bigl[\mathbf 1+s\,a_1+s^2 a_2+\cdots\bigr]$, and the integrated fourth
Seeley--DeWitt coefficient is given by Gilkey's formula \cite{Gilkey}
\begin{equation}\label{eq:gilkey}
\int a_4=\frac{1}{(4\pi)^2\,360}\int\!\sqrt g\;\mathrm{tr}\Bigl\{60RE+180E^2+30\,\Omega_{\mu\nu}\Omega^{\mu\nu}
+\mathbf 1\bigl(5R^2-2R_{\mu\nu}^2+2R_{\mu\nu\rho\sigma}^2\bigr)\Bigr\},
\end{equation}
where $\Omega_{\mu\nu}$ is the curvature of the connection on the bundle. The conformal anomaly of
a single field equals $\zeta(0)$, which is exactly the integrated $a_4$; on the unit $S^4$ one has
the Einstein-space identities $R_{\mu\nu}^2=R^2/4$, $R_{\mu\nu\rho\sigma}^2=R^2/6$, $W^2=0$, and
$\int E_4=64\pi^2$ (Theorem~\ref{thm:gbc}), so only the Euler part survives and a single field
contributes $\zeta(0)=-4a$ where $a$ is its type-A coefficient. We evaluate three operators.

\smallskip
\emph{(i) Conformally coupled scalar}, $D=-\Box+\tfrac R6$ ($E=-\tfrac R6$, $\dim 1$,
$\Omega=0$). The bracket in \eqref{eq:gilkey} reduces to $-\tfrac16 R^2$ and yields
$\zeta(0)_{\mathrm{scalar}}=-\tfrac1{90}$, i.e. $a_{\mathrm{scalar}}=\tfrac1{360}$ --- the
textbook value. This serves as a normalization check on the whole computation.

\emph{(ii) Hodge--de~Rham Laplacian on $1$-forms}, with Weitzenb\"ock potential $E=-\mathrm{Ric}$
(so $\mathrm{tr}\,E=-R$, $\mathrm{tr}\,E^2=R^2/4$, $\dim 4$,
$\mathrm{tr}\,\Omega_{\mu\nu}\Omega^{\mu\nu}=-R_{\mu\nu\rho\sigma}^2=-R^2/6$). Substituting,
\[
\zeta(0)_{\Delta_1}=\int\avg{T}_{\Delta_1}=-\tfrac{2}{45}.
\]

\emph{(iii) Minimal scalar ghost}, $E=0$, $\dim 1$: $\zeta(0)_{\Delta_0}=\int\avg{T}_{\Delta_0}=\tfrac{29}{90}$.

\smallskip
The photon partition function is gauge-fixed with two Faddeev--Popov scalar ghosts removing the
gauge and longitudinal modes, $\ln Z=-\tfrac12\ln\det\Delta_1+\ln\det\Delta_0^{\mathrm{ghost}}$, so
its anomaly is the operator combination $\Delta_1-2\Delta_0$:
\begin{equation}\label{eq:zetacombo}
\zeta(0;\mathrm{Maxwell};S^4)=\zeta(0)_{\Delta_1}-2\,\zeta(0)_{\Delta_0}
=-\tfrac{2}{45}-2\cdot\tfrac{29}{90}=-\tfrac{4}{90}-\tfrac{58}{90}=-\tfrac{62}{90}=-\tfrac{31}{45}.
\end{equation}
Since a single field gives $\zeta(0)=-4a$, we read off
$\ag=-\tfrac14\bigl(-\tfrac{31}{45}\bigr)=\dfrac{31}{180}$.
\end{proof}

\begin{remark}
$\ag$ is thus \emph{computed}, not adopted from a citation. The recovered scalar value $1/360$ is
an independent check of the normalization of \eqref{eq:gilkey}.
\end{remark}

\section{The topological prefactor: Gauss--Bonnet--Chern}\label{sec:gbc}
\begin{theorem}[Gauss--Bonnet--Chern; the prefactor $4=2\chi(S^4)$]\label{thm:gbc}
For a closed oriented Riemannian four-manifold $M$,
$\displaystyle\int_M E_4\,\sqrt g\,d^4x=32\pi^2\,\chi(M)$, where $\chi$ is the Euler
characteristic. On $S^4$, $\chi=2$, so $\int_{S^4}E_4=64\pi^2$, and with the conformal
normalization $16\pi^2$ in \eqref{eq:anomaly},
\[
\int_{S^4}\avg{T^\mu{}_\mu}=\frac{\ag}{16\pi^2}\int_{S^4}E_4=\frac{\ag}{16\pi^2}\cdot 64\pi^2=4\ag,
\qquad
\frac{\int_{S^4}E_4}{16\pi^2}=4=2\,\chi(S^4).
\]
The prefactor $4$ is the Atiyah--Singer index for the Euler characteristic; it is geometric and
not a free parameter. \cite{Chern,AtiyahSinger}
\end{theorem}

\section{The hemisphere curvature charge $\CQ=8\pi^2$}\label{sec:Q}
\begin{definition}[$Q$-curvature, Paneitz operator]
On a four-manifold the Branson $Q$-curvature is the local scalar
$Q_4=-\tfrac16\Box R-\tfrac12 R_{\mu\nu}R^{\mu\nu}+\tfrac16 R^2$ (up to convention), the natural
fourth-order analogue of the Gauss curvature; it transforms under Weyl rescaling by the
fourth-order Paneitz operator $P_4$. On an Einstein manifold $\Box R=0$ and $Q_4$ reduces to
$\tfrac16(R^2-3R_{\mu\nu}R^{\mu\nu})$.
\end{definition}

\begin{lemma}[Hemisphere $Q$-charge]\label{lem:qcharge}
On the round unit hemisphere $D^4=\{0\le\theta\le\pi/2\}\subset S^4$,
$\displaystyle\CQ\equiv\int_{D^4}Q_4\,\sqrt g\,d^4x=8\pi^2$.
\end{lemma}
\begin{proof}
On the unit $S^4$ (Einstein) $R=12$, $R_{\mu\nu}=3g_{\mu\nu}$ so $R_{\mu\nu}R^{\mu\nu}=36$, and
$\Box R=0$; hence $Q_4=\tfrac16(R^2-3R_{\mu\nu}R^{\mu\nu})=\tfrac16(144-108)=6$. With volume
element $\sqrt g=\sin^3\theta$, $\mathrm{Vol}(S^3)=2\pi^2$, and
\[
\int_0^{\pi/2}\sin^3\theta\,d\theta=\Bigl[-\cos\theta+\tfrac13\cos^3\theta\Bigr]_0^{\pi/2}=\tfrac23,
\]
we obtain $\int_{D^4}Q_4\,\sqrt g\,d^4x=6\cdot 2\pi^2\cdot\tfrac23=8\pi^2.$
\end{proof}
\begin{remark}
$\CQ=8\pi^2$ is exact and scale-free. The ratio $\int_{D^4}Q_4/\int_{D^4}E_4=1/4$ (result $R12$
below) fixes the relative $Q$-vs-$E$ normalization used throughout.
\end{remark}

\section{Spectral backbone: the twelve results and the BFK check}\label{sec:spectral}
The construction rests on twelve $\zeta$-function results on $S^4$, its hemisphere $D^4$, and the
equatorial $S^3$. They are stated compactly; $R1$ is the load-bearing one and is derived twice.

\begin{description}[leftmargin=2.6em,nosep,itemsep=2pt]
\item[R1.] $\zeta(0;\mathrm{Maxwell};S^4)=-4\ag=-31/45$. Two derivations agree: the integrated
trace-anomaly route ($-\ag\int E_4/16\pi^2=-4\ag$) and the spectral combination
\eqref{eq:zetacombo}. Each unit of Euler characteristic contributes $-2\ag$; $\chi(S^4)=2$.
\item[R2.] $\zeta(0;\mathrm{DtN};\mathrm{scalar};S^3)=-1$. The scalar Dirichlet-to-Neumann operator
on the equatorial $S^3$ (harmonic extension into the cap) has spectrum
$\lambda_\ell=\ell(\ell+2)/(\ell+1)$, multiplicity $(\ell+1)^2$; the analytic continuation gives
$\zeta(0)=\zeta_R(-2)-1=-1$.
\item[R3.] $\zeta(0;\mathrm{DtN};\text{1-form};S^3)=1$.
\item[R4.] $\zeta(0;\mathrm{DtN};\mathrm{Maxwell};S^3)=3=\dim S^3$ (gauge-fixed Maxwell DtN).
\item[R5.] $\zeta(0;\mathrm{Maxwell};S^3)=1$ (photon equatorial identity).
\item[R6.] BFK identity: with $R1$ and the two bookkeepings below,
$2\zeta(0;D^4)-\zeta(0;\mathrm{DtN})=-31/45$.
\item[R7.] Photon Casimir energy on $S^3$: $E_{\mathrm{Cas}}=11/120$.
\item[R8.] $\det{}'\Lambda^{\mathrm{Maxwell}}_{S^3}=1/(2\pi^3)$.
\item[R9.] $\eta$ of the transverse $1$-form DtN equals its $\zeta(0)=1$, conformally invariant
(distinct from the de~Rham APS $\eta_{\mathrm{APS}}(S^3)=0$ used in the topological quantization).
\item[R10.] A Robin condition $(\partial_n+\alpha)u=0$ at the equator shifts the DtN spectrum by
$\alpha=\tan\ag\approx\ag$ --- the analytic origin of the cap-saddle drive (\S\ref{sec:cap}).
\item[R11.] Under the Robin shift, $\zeta(0;\Lambda_{\mathrm{Robin}})$ is linear in $\alpha$.
\item[R12.] $\int_{D^4}Q_4/\int_{D^4}E_4=1/4$ exactly.
\end{description}

\paragraph{Two consistent BFK bookkeepings.} The Burghelea--Friedlander--Kappeler gluing identity
relates the $\zeta$-determinant of an operator on a closed manifold to those on two pieces plus a
Dirichlet-to-Neumann boundary term. Splitting $S^4$ into two hemispheres, the closed value $R1$ is
recovered in two internally consistent ways differing only in the gauge/ghost assignment:
\begin{align}
\text{(a) scalar-DtN split:}\quad
&\zeta(0;D^4)=-2\ag-\tfrac12=-\tfrac{38}{45},
&&2\bigl(-\tfrac{38}{45}\bigr)-(-1)=-\tfrac{31}{45}\ \checkmark\\
\text{(b) gauge-fixed split:}\quad
&\zeta(0;D^4)=\tfrac{52}{45},
&&2\bigl(\tfrac{52}{45}\bigr)-3=-\tfrac{31}{45}\ \checkmark
\end{align}
Their agreement is an internal check on the spectral arithmetic; $R1$ itself is independently
\emph{derived} from the heat kernel in \S\ref{sec:heat}.

\section{The hierarchy face: sequestering on the hemisphere}\label{sec:seq}
A naive vacuum energy scales as $\Mpl^4$. The Kaloper--Padilla \emph{sequestering} mechanism
renders the cosmological constant a global Lagrange multiplier rather than a local source: a global
constraint removes the radiatively unstable bulk contribution, leaving only the boundary/topological
piece. On the hemisphere this leaves precisely the anomaly term: with $\int_{D^4}=\tfrac12\int_{S^4}$,
\begin{equation}\label{eq:sequester}
\tfrac12\,\Lambda_{\mathrm{eff}}\,V_4(D^4)=\ag .
\end{equation}
The point we use is structural: after sequestering, the residual cosmological term is fixed by the
\emph{dimensionless} anomaly charge, not by the cutoff. This is what makes a parameter-free ratio
possible; it addresses the hierarchy face of the problem (\S\ref{sec:problem}, (i)).

\section{The infrared lift by coherent spectral accumulation}\label{sec:lift}
The local anomaly density on de~Sitter space is
$\rho_{\mathrm{anom}}=\tfrac14\avg{T^\mu{}_\mu}=\tfrac{\ag}{8\pi^2}\,3H^4$, giving
$\rho_{\mathrm{anom}}/\rho_{\mathrm{crit}}=\tfrac{\ag}{8\pi^2}(H/\Mpl)^2\sim10^{-122}$ --- a local
reading alone is far too small. The closure of this gap is a property of the spectrum, not an added
assumption.

\begin{proposition}[Coherent accumulation]\label{prop:lift}
The gauge-fixed Maxwell operator on $S^4$ has discrete spectrum $\lambda_n=n(n+2)H^2$ with
multiplicity $d_n=2n(n+2)$, $n\ge1$ \cite{Allen,Higuchi}. These modes are fixed and phase-locked by
the global geometry, so the $\zeta$-regularized vacuum sum is \emph{coherent} (linear in the mode
count $N$) rather than stochastic ($\sqrt N$). The bulk count to the cutoff $n_{\max}=\Mpl/H$ is
\[
N_{\mathrm{tot}}=\sum_{n=1}^{n_{\max}}2n(n+2)=\tfrac23 n_{\max}^3+3n_{\max}^2+\tfrac73 n_{\max}
\sim\tfrac23\Bigl(\frac{\Mpl}{H}\Bigr)^3,
\]
and the horizon-area projection reduces this to
$N_{\mathrm{eff}}\sim(\Mpl/H)^2=S_{dS}/\pi$, the Gibbons--Hawking de~Sitter entropy. The
$(\Mpl/H)^2$ scaling is exact; its $O(1)$ prefactor is convention-dependent.
\end{proposition}

\begin{corollary}
$N_{\mathrm{eff}}\,\rho_{\mathrm{anom}}/\rho_{\mathrm{crit}}=\ag/8\pi^2$, i.e. the dimensionless
ratio of \S\ref{sec:friedmann}. The magnitude of the \emph{ratio} --- the quantity actually
measured --- is thus fixed without a free parameter.
\end{corollary}

\section{The integrated invariant $4\ag$ by two independent routes}\label{sec:invariant}
\begin{theorem}[Spectral route]\label{thm:zeta}
$\displaystyle\int_{S^4}\avg{T^\mu{}_\mu}=4\ag=\bigl\lvert\zeta(0;\mathrm{Maxwell};S^4)\bigr\rvert.$
\end{theorem}
\begin{proof}
The conformal anomaly is the constant-Weyl variation of the one-loop effective action,
$\delta_\sigma\Gamma=-\tfrac12\zeta(0)\delta\sigma$, and $\zeta(0)$ is the integrated $a_4$. By
$R1$, $\zeta(0;\mathrm{Maxwell};S^4)=-31/45=-4\ag$.
\end{proof}

\begin{theorem}[Scale-flow route]\label{thm:scaleflow}
Under $g\to e^{2\sigma}g$ the anomaly-induced (Riegert--Mottola) action satisfies the
Wess--Zumino consistency relation with
$\displaystyle\frac{dS_{\mathrm{anom}}}{d\sigma}=\frac{\ag}{16\pi^2}\int_{S^4}E_4=4\ag,$
again because $\int_{S^4}E_4=64\pi^2$. \cite{Riegert1984,Mottola}
\end{theorem}

\begin{remark}
Theorems~\ref{thm:zeta} and~\ref{thm:scaleflow} fix the number $4\ag$ by two constructions sharing
no intermediate step and no free $O(1)$ factor. The value $4\ag$ is therefore prefactor-independent.
\end{remark}

\section{The variational mechanism, stability, and the sign}\label{sec:cap}
\begin{proposition}[Cap-saddle and the open universe]\label{prop:sign}
Result $R10$ shows that a Robin condition at the equator shifts the DtN spectrum by
$\alpha=\tan\ag$. Extremizing the resulting boundary effective action over the equatorial position
$\delta\theta$ gives the stationarity equation
\[
\ag\cos^4(\delta\theta)=\sin(\delta\theta),
\]
with a stable solution $\delta\theta^*\approx0.164>0$. Stability follows from the positive discrete
spectrum of the Laplacian on the compact $S^4$ (a maximum of the Euclidean action is a minimum of
the energy after Wick rotation). The positive sign of $\delta\theta^*$ forces $\Ok>0$: an
\emph{open} universe. The charge $\CQ=8\pi^2$ is always evaluated on the standard hemisphere, so the
leading ratio $\lvert\Ok\rvert/\OL=\ag/8\pi^2$ carries no $\delta\theta^*$ correction; the cap-saddle
enters only at the refinement level.
\end{proposition}

\section{Uniqueness: why the photon}\label{sec:uniq}
We record why no other field replaces the photon and why no alternative mechanism closes the
absolute value. \emph{(a) Field content:} massive fields decouple as $m^2/\Mpl^2$; among massless
fields the graviton's type-A coefficient enters only with propagating quantum gravity, which we do
not invoke, while neutrinos (if massless) and gluons are either not conformally coupled in the
relevant sense or are confined; the unbroken $U(1)$ photon is the clean contributor. \emph{(b)
Mechanism:} candidate dynamical closures of the absolute value (a scalaron with a
Coleman--Weinberg potential; identifying a second operator with dark matter; a stochastic rather
than coherent sum) were each examined and found either to give the wrong number
($1/(12\ag)\approx0.48$, not $0.69$), to be locally negligible, or to violate the coherence shown
in Proposition~\ref{prop:lift}. The surviving statement is the proven \emph{ratio} together with the
single identifying postulate of \S\ref{sec:postulate}.

\section{The parameter-free ratio and the Friedmann translation}\label{sec:friedmann}
\begin{theorem}[Central relation]\label{thm:central}
With $\CQ=8\pi^2$ (Lemma~\ref{lem:qcharge}) and $\ag=31/180$ (Theorem~\ref{thm:agamma}),
\begin{equation}\label{eq:central}
\frac{\lvert\Ok\rvert}{\OL}=\frac{\ag}{\CQ}=\frac{\ag}{8\pi^2}=\frac{31}{1440\pi^2}\approx2.181\times10^{-3}.
\end{equation}
The relation contains no scale, no Newton constant, no Planck mass.
\end{theorem}

\paragraph{Eight-step assembly.}
\emph{(1)} $\int_{S^4}\avg{T}=4\ag$ (Thm.~\ref{thm:gbc}); \emph{(2)} sequestering on $D^4$,
Eq.~\eqref{eq:sequester}; \emph{(3)} $\CQ=8\pi^2$ (Lem.~\ref{lem:qcharge}); \emph{(4)} Wess--Zumino
linear response (Thm.~\ref{thm:scaleflow}); \emph{(5)} the dimensionless quotient
\eqref{eq:central}; \emph{(6)} photon selection (\S\ref{sec:uniq}); \emph{(7)} the Euclidean
cap-saddle (Prop.~\ref{prop:sign}) and its Wick rotation to a Lorentzian Friedmann patch;
\emph{(8)} the Friedmann translation below. Steps $1$--$6$ and the Euclidean sub-steps of $7$ are
theorem-level; the one genuinely chosen element is the present-epoch identification isolated in
\S\ref{sec:postulate}.

\paragraph{Friedmann translation.} Writing $\overline\Lambda=8\pi^2/\ag=1440\pi^2/31\approx458.4$,
the postulate $\OL=4\ag$ (\S\ref{sec:postulate}) together with $\OL=\overline\Lambda\lvert\Ok\rvert$
fixes
\begin{equation}\label{eq:OK}
\lvert\Ok\rvert=\frac{\OL}{\overline\Lambda}=\frac{\ag^2}{2\pi^2}\approx1.50\times10^{-3}
\quad(\text{open, }\Ok>0),
\end{equation}
and the flatness budget $\Om+\OL+\Ok=1$ then \emph{predicts}
$\Om=1-\OL-\lvert\Ok\rvert\approx0.3096$, consistent with the Planck~2018 value
$\Om=0.3111\pm0.0056$; equivalently $\OL=0.6889$ matches Planck to four decimals
\cite{Planck2018}.

\section{The single physical postulate}\label{sec:postulate}
Everything in \S\S\ref{sec:anomaly}--\ref{sec:friedmann} is Euclidean and geometric. One step
connects the proven Euclidean ratio to the present-epoch Friedmann observable.

\begin{postulate}[Anomaly--curvature correspondence, ``A5/P5'']\label{post:A5}
The conformal-anomaly ratio $\ag/\CQ$ computed on the Euclidean hemisphere, carried through the
Wick rotation to a Lorentzian Friedmann patch, equals the observed present-epoch spatial-curvature
parameter; equivalently, the dark-energy density saturates the integrated invariant, $\OL=4\ag$.
\end{postulate}

\begin{remark}[Exact status of Postulate~\ref{post:A5}]
The Euclidean half --- the topological identity $\Ok^{\mathrm{conformal}}=\ag/\CQ$ --- is a
\emph{theorem}: Cauchy's functional equation gives the linearity, a Chern--Weil basis change gives
the $1/8\pi^2$, and the Atiyah--Patodi--Singer index theorem on the de~Rham complex gives the
integer normalization (up to a bounded boundary refinement $\lvert\epsilon_\partial\rvert\le1.5\%$).
The Wick-rotation kinematics are \emph{derivable} in the saddle-point approximation. What remains a
genuine \emph{physical hypothesis} is the identification of the resulting Euclidean quantity with
the present-epoch observable under the canonical $\Lambda$-dominated reference epoch. It is
physically motivated --- $4\ag$ is simultaneously the spectral $\zeta$-invariant
(Thm.~\ref{thm:zeta}) and the Weyl scale-anomaly coefficient (Thm.~\ref{thm:scaleflow}) --- but is
\emph{not} forced dynamically: the anomaly action $S_{\mathrm{anom}}(\sigma)=4\ag\,\sigma$ is linear
in $\sigma$ and has no stationary point, and the dynamical scalaron route returns $1/(12\ag)\approx
0.48$ rather than $0.69$. We therefore state it as a founding postulate, of the same logical type as
the postulated principles of the established physical theories (the equivalence principle, the
Schr\"odinger equation, the Born rule, the gauge principle), with the distinguishing feature that
its value is independently fixed by the two invariants of \S\ref{sec:invariant}. It is not claimed
to be a derived theorem, and it is the proper target of any falsification attempt.
\end{remark}

\section{Predictions, observational status, forecasts}\label{sec:obs}
Given Postulate~\ref{post:A5}:
\[
\OL=4\ag\approx0.6889,\qquad \Ok\simeq+1.5\times10^{-3}\ (\text{open}),\qquad w=-1 .
\]
The ratio $\lvert\Ok\rvert/\OL=\ag/8\pi^2$ of Theorem~\ref{thm:central} holds \emph{independently}
of the postulate. All are falsifiable: a measured $w$ differing from $-1$ at high significance, a
closed universe ($\Ok<0$), or a curvature/dark-energy ratio departing from $2.181\times10^{-3}$
each refutes the framework. The equation-of-state value $w=-1$ is presently in $\sim3$--$4\sigma$
tension with DESI~DR2; DESI~DR3 and Euclid will decide. This is a sharp, risky prediction, not a
retrodiction.

\section{Logical status of every claim}\label{sec:ledger}
\begin{center}
\begin{tabular}{lll}
\toprule
Claim & Status & Method / reference \\
\midrule
$\ag=31/180$ & proved & Gilkey $a_4$, three operators (Thm.~\ref{thm:agamma}) \\
$\int_{S^4}E_4=64\pi^2$; prefactor $4=2\chi$ & proved & Gauss--Bonnet--Chern (Thm.~\ref{thm:gbc}) \\
$\CQ=8\pi^2$ & proved & elementary integral (Lem.~\ref{lem:qcharge}) \\
spectral backbone $R1$--$R12$ & proved & \S\ref{sec:spectral}; $R1$ twice (heat kernel + BFK) \\
$\int_{S^4}\avg T=4\ag$ & proved (two routes) & $\zeta$-determinant; scale-flow (Thms.~\ref{thm:zeta},~\ref{thm:scaleflow}) \\
ratio $\ag/8\pi^2$ & proved, parameter-free & Thm.~\ref{thm:central} \\
sequestering suppression & literature mechanism & Kaloper--Padilla (\S\ref{sec:seq}) \\
lift $(\Mpl/H)^2$ scaling & proved (scaling) & coherent spectral sum (Prop.~\ref{prop:lift}) \\
sign $\Ok>0$ & proved & spectral positivity (Prop.~\ref{prop:sign}) \\
photon uniqueness & argued (exhaustion) & \S\ref{sec:uniq} \\
$\OL=4\ag$ (absolute value) & \textbf{postulated} & identification A5/P5 (Post.~\ref{post:A5}) \\
numerical $\Om,\OL,\Ok$; $w=-1$ & empirical & DESI~DR3 / Euclid (open) \\
\bottomrule
\end{tabular}
\end{center}

\noindent
Every ``proved'' row is offered for independent verification on its own terms. The single
load-bearing hypothesis is the row marked \textbf{postulated}; any falsification of the framework
should be directed there or at the empirical predictions of \S\ref{sec:obs}.

\begin{thebibliography}{99}
\bibitem{Gilkey} P.~B.~Gilkey, \emph{Invariance Theory, the Heat Equation, and the
Atiyah--Singer Index Theorem}, 2nd ed., CRC Press (1995).
\bibitem{Birrell1982} N.~D.~Birrell, P.~C.~W.~Davies, \emph{Quantum Fields in Curved Space},
Cambridge (1982).
\bibitem{ChristensenDuff1980} S.~M.~Christensen, M.~J.~Duff, Nucl. Phys. B \textbf{170} (1980) 480.
\bibitem{Vassilevich2003} D.~V.~Vassilevich, \emph{Heat kernel expansion: user's manual}, Phys.
Rep. \textbf{388} (2003) 279.
\bibitem{Duff1994} M.~J.~Duff, \emph{Twenty years of the Weyl anomaly}, Class. Quantum Grav.
\textbf{11} (1994) 1387.
\bibitem{DeserSchwimmer} S.~Deser, A.~Schwimmer, Phys. Lett. B \textbf{309} (1993) 279.
\bibitem{Chern} S.-S.~Chern, Ann. Math. \textbf{45} (1944) 747.
\bibitem{AtiyahSinger} M.~F.~Atiyah, I.~M.~Singer, Ann. Math. \textbf{87} (1968) 484.
\bibitem{Riegert1984} R.~J.~Riegert, Phys. Lett. B \textbf{134} (1984) 56.
\bibitem{Mottola} E.~Mottola, anomaly-induced effective action (and references therein).
\bibitem{Allen} B.~Allen, vector two-point functions on maximally symmetric spaces (1986).
\bibitem{Higuchi} A.~Higuchi, J. Math. Phys. \textbf{28} (1987) 1553.
\bibitem{BGT1995} M.~Bucher, A.~Goldhaber, N.~Turok, Phys. Rev. D \textbf{52} (1995) 3314.
\bibitem{HawkingTurok1998} S.~W.~Hawking, N.~Turok, Phys. Lett. B \textbf{425} (1998) 25.
\bibitem{Planck2018} Planck Collaboration, \emph{Planck 2018 results. VI}, A\&A \textbf{641} (2020) A6.
\end{thebibliography}

\end{document}
